\chapter{Estado del arte}

La investigación sobre locomoción cuadrúpeda combina modelado, simulación y validación experimental. Plataformas como MIT Cheetah~3 y ANYmal han servido para evaluar arquitecturas de control capaces de coordinar múltiples articulaciones y mantener la estabilidad durante desplazamientos dinámicos~\cite{bledt2018cheetah3,hutter2016anymal}. Estos trabajos evidencian la importancia de validar el controlador tanto en condiciones nominales como ante perturbaciones.

Una línea de trabajo utiliza optimización para generar movimientos de cuerpo completo. Bellicoso \textit{et al.} formulan la locomoción como un problema de optimización no lineal en línea, mientras que Neunert \textit{et al.} incorporan la dinámica y los contactos en un esquema de control predictivo no lineal~\cite{bellicoso2018dynamic,neunert2018wholebody}. En ambos casos, las restricciones físicas del robot forman parte explícita del problema de control.

Otra aproximación emplea modelos simplificados para reducir el costo computacional. El controlador predictivo convexo desarrollado para MIT Cheetah~3 calcula fuerzas de reacción y referencias de movimiento en tiempo real, permitiendo ejecutar locomoción dinámica y responder a perturbaciones externas~\cite{dicarlo2018mpc}. Para un prototipo educativo, esta línea ofrece criterios útiles para separar la generación nominal de marcha de la capa de estabilización.

La planificación tradicional de marcha continúa siendo relevante, especialmente cuando se busca una implementación interpretable y segura. Los modelos cinemáticos y dinámicos permiten definir secuencias de apoyo, verificar el espacio alcanzable y conservar condiciones de estabilidad durante marchas lentas~\cite{zamani2011dynamics}. Una máquina de estados finitos puede utilizar estas secuencias como referencia nominal antes de incorporar estrategias adaptativas.

El aprendizaje por refuerzo constituye una línea complementaria. Se ha observado que políticas entrenadas con objetivos relacionados con locomoción y consumo energético pueden producir comportamientos coordinados y diferentes patrones de marcha~\cite{heess2017emergence,fu2021energy}. Sin embargo, la transferencia al robot físico exige controlar el dominio de entrenamiento, las acciones permitidas y los mecanismos de seguridad.

La simulación dinámica permite repetir experimentos y modificar parámetros sin exponer el hardware. MuJoCo fue diseñado para simulación física y control basado en modelos, y se utiliza como soporte para desarrollar y evaluar estrategias antes de transferirlas al robot~\cite{todorov2012mujoco,tassa2018controlsuite}. En esta tesis se combinará la simulación con ROS~2 para conservar una arquitectura modular de comunicación y ejecución.

\section{Comparación con desarrollos actuales}

Los sistemas de referencia emplean actuadores, sensores y computadores de alto desempeño~\cite{bledt2018cheetah3,hutter2016anymal}. El presente proyecto trabaja con una plataforma de laboratorio de recursos más limitados; por ello, propone una marcha discreta como línea base y una política aprendida como capa correctiva limitada. Esta separación busca conservar interpretabilidad y facilitar una transferencia progresiva y segura.

\section{Vacío identificado}

Los trabajos revisados demuestran resultados avanzados en plataformas especialmente diseñadas para locomoción dinámica. Existe una oportunidad de evaluar qué parte de esas estrategias puede trasladarse a un robot educativo de bajo costo, con actuadores de posición y capacidad computacional limitada. El aporte de esta tesis será una implementación documentada y una comparación experimental bajo métricas comunes, no la formulación de un nuevo algoritmo de aprendizaje.

\section{Punto de partida del proyecto}

El proyecto partió de una plataforma Nova Spot Micro sin una línea experimental
documentada. Sobre esta base se desarrollaron los modelos nominales, el
controlador discreto, la instrumentación simulada y el protocolo experimental.
La arquitectura prevista para aprendizaje por refuerzo se delimitó como una
corrección de referencias nominales, con límites de acción y un supervisor
independiente; su evaluación comparativa permanece pendiente.
